% ===== PRÉAMBULE CANONIQUE — à coller VERBATIM en tête de CHAQUE .tex =====
\documentclass[11pt,a4paper]{article}
\usepackage[utf8]{inputenc}\usepackage[T1]{fontenc}\usepackage{lmodern}
\usepackage[french]{babel}
\usepackage{amsmath,amssymb,mathtools}\usepackage{icomma}
\usepackage{siunitx}\sisetup{locale=FR}  % \qty{}{}, \si{}, \ang{}
\usepackage{xcolor}\usepackage{colortbl}
\usepackage{tikz}\usepackage{pgfplots}\pgfplotsset{compat=1.18}
\usepackage[siunitx,europeanresistors]{circuitikz} % circuits (élec.)
\usepackage[most]{tcolorbox}\usepackage{enumitem}\usepackage{array,booktabs}
\usepackage[a4paper,margin=2.2cm]{geometry}
\usetikzlibrary{arrows.meta,calc,angles,quotes,positioning,patterns,%
  backgrounds,shapes.geometric,decorations.pathmorphing,decorations.markings,intersections}
\definecolor{primary}{RGB}{222,49,99}\definecolor{primarydark}{RGB}{150,22,55}
\definecolor{defcol}{RGB}{0,114,178}\definecolor{methcol}{RGB}{52,92,148}
\definecolor{excol}{RGB}{0,135,90}\definecolor{attncol}{RGB}{200,40,55}
\tcbset{boxbase/.style={enhanced,breakable,boxrule=0.9pt,arc=2.5pt,
  left=3mm,right=3mm,top=2.3mm,bottom=2.3mm,fonttitle=\bfseries,coltitle=white}}
% --- boîtes TUTORIEL (noms EXACTS, ne pas renommer) ---
\newtcolorbox{cle}[1][]{boxbase,colback=defcol!6,colframe=defcol,
  title={\textsc{À retenir}\ifx\relax#1\relax\relax\else\textmd{ : \itshape #1}\fi}}
\newtcolorbox{methode}[1][]{boxbase,colback=methcol!7,colframe=methcol,sharp corners,
  title={\textsc{Méthode}\ifx\relax#1\relax\relax\else\textmd{ --- \itshape #1}\fi}}
\newtcolorbox{exemplebox}[1][]{boxbase,colback=excol!5,colframe=excol,
  title={\textsc{Exemple}\ifx\relax#1\relax\relax\else\textmd{ : \itshape #1}\fi}}
\newtcolorbox{piege}[1][]{boxbase,colback=attncol!6,colframe=attncol,
  title={\textsc{Piège}\ifx\relax#1\relax\relax\else\textmd{ : \itshape #1}\fi}}
\newtcolorbox{remarque}[1][]{boxbase,colback=black!4,colframe=black!45,
  title={\textsc{Remarque}\ifx\relax#1\relax\relax\else\textmd{ : \itshape #1}\fi}}
% --- boîtes EXERCICE / CORRIGÉ (noms EXACTS, auto-comptées) ---
\newtcolorbox[auto counter]{exo}[1][]{boxbase,colback=primary!4,colframe=primary,
  title={\textsc{Exercice}~\thetcbcounter\ifx\relax#1\relax\relax\else\textmd{ --- \itshape #1}\fi}}
\newtcolorbox[auto counter]{corrige}[1][]{boxbase,colback=excol!4,colframe=excol,
  title={\textsc{Corrigé}~\thetcbcounter\ifx\relax#1\relax\relax\else\textmd{ --- \itshape #1}\fi}}
\newcommand{\vv}[1]{\vec{#1}}\newcommand{\ds}{\displaystyle}
\newcommand{\R}{\mathbb{R}}\newcommand{\N}{\mathbb{N}}\newcommand{\Z}{\mathbb{Z}}
% (un \usepackage{fancyhdr}+en-tête est possible ; il est ignoré côté web)
% =========================================================================
\begin{document}
\sisetup{per-mode=symbol}

\begin{exo}[distance de freinage]
Une voiture roule à \qty{108}{\kilo\metre\per\hour} et freine avec une décélération constante $a=\qty{-5}{\metre\per\second\squared}$.
\begin{enumerate}[label=\alph*)]
  \item Convertis la vitesse en \si{\metre\per\second}.
  \item Quelle distance parcourt-elle avant de s'arrêter ?
  \item Combien de temps dure le freinage ?
\end{enumerate}
\end{exo}

\begin{exo}[du repos à la vitesse de croisière]
Une voiture partant du repos atteint \qty{72}{\kilo\metre\per\hour} sur une distance de \qty{100}{\metre}
en accélération uniforme.
\begin{enumerate}[label=\alph*)]
  \item Quelle est son accélération ?
  \item En combien de temps atteint-elle cette vitesse ?
\end{enumerate}
\end{exo}

\begin{exo}[deux phases : accélération puis MRU]
Un métro démarre avec $a=\qty{2}{\metre\per\second\squared}$ jusqu'à \qty{20}{\metre\per\second}, puis roule à
vitesse constante pendant \qty{5}{\second}.
\begin{enumerate}[label=\alph*)]
  \item Combien de temps dure la phase d'accélération et quelle distance couvre-t-elle ?
  \item Quelle distance parcourt-il pendant la phase à vitesse constante ?
  \item Quelle est la distance totale parcourue ?
\end{enumerate}
\end{exo}

\begin{exo}[glissade sur un plan incliné]
Un bloc part du repos et glisse sans frottement le long d'un plan incliné à \ang{30}
(l'accélération le long de la pente vaut $a=g\sin\ang{30}$, avec $g=\qty{10}{\metre\per\second\squared}$).
\begin{center}
\begin{tikzpicture}[scale=0.9,>=Stealth]
  \draw[thick] (0,0)--(5,0);
  \draw[thick] (0,0)--(4.33,2.5);
  \draw (0.8,0) arc (0:30:0.8); \node at (1.15,0.22){\small\ang{30}};
  \fill[primary!70] (2.6,1.5) rectangle ++(0.45,0.32);
  \draw[->,line width=1pt,primary] (3.0,1.75)--(3.9,1.23) node[right]{$a$};
\end{tikzpicture}
\end{center}
\begin{enumerate}[label=\alph*)]
  \item Quelle est l'accélération du bloc ?
  \item Quelle vitesse atteint-il après avoir parcouru \qty{10}{\metre} le long de la pente ?
\end{enumerate}
\end{exo}

\begin{exo}[lire les paramètres dans l'équation horaire]
La position d'un mobile est $x(t)=1{,}5\,t^2+4t+2$ (en \si{\metre}, $t$ en \si{\second}).
\begin{enumerate}[label=\alph*)]
  \item Identifie $x_0$, $v_0$ et l'accélération $a$.
  \item Le mouvement est-il accéléré ou décéléré ? Quelle est la vitesse à $t=\qty{2}{\second}$ ?
\end{enumerate}
\end{exo}

\end{document}
